\documentclass[11pt]{article}

% ---------- Layout ----------
\usepackage[margin=0.8in, top=0.6in, bottom=0.7in, headsep=10pt]{geometry}
\usepackage{times}
\usepackage[skip=3pt]{parskip}
\usepackage{titlesec}
\usepackage{enumitem}
\usepackage{fancyhdr}
\usepackage{hyperref}
\usepackage{graphicx}

\hypersetup{
    colorlinks=false,
    pdfborder={0 0 0},
    pdftitle={Week 2 Progress Report},
    pdfauthor={Aryan Thakkar, Nikkesh Parekh, Hitaksh Thaker, Maitri Parekh}
}

% ---------- Section styling ----------
\titleformat{\section}
  {\normalfont\large\bfseries}
  {\thesection}{0.6em}{}

\titlespacing*{\section}{0pt}{6.5pt}{2.5pt}

\renewcommand{\thesection}{\arabic{section}}

% ---------- List spacing ----------
\setlist[itemize]{
    itemsep=1pt,
    topsep=2pt,
    parsep=0pt,
    leftmargin=1.4em
}

\setlist[enumerate]{
    itemsep=1pt,
    topsep=2pt,
    parsep=0pt,
    leftmargin=1.6em
}

% ---------- Header / footer ----------
\pagestyle{fancy}
\fancyhf{}

\renewcommand{\headrulewidth}{0.4pt}
\renewcommand{\footrulewidth}{0.4pt}

\fancyhead[L]{\small\textsc{ECE 501 --- Digital Image Processing}}
\fancyhead[R]{\small Week 2 Progress Report}

\fancyfoot[L]{\small\textit{Galaxy Morphology Classification --- Week 2}}
\fancyfoot[R]{\small\thepage}

% ============================================================
% DOCUMENT
% ============================================================

\begin{document}

\thispagestyle{fancy}

% ============================================================
% TITLE BLOCK
% ============================================================

\begin{center}

\Large\textbf{Morphological Classification of Galaxies in HST and JWST Images Using Classical DIP}

\vspace{3pt}

\normalsize\textbf{Weekly Progress Report --- Week 2}

\vspace{4pt}

\small
\textbf{Aryan Thakkar}
\quad $\cdot$ \quad
\textbf{Nikkesh Parekh}
\quad $\cdot$ \quad
\textbf{Hitaksh Thaker}
\quad $\cdot$ \quad
\textbf{Maitri Parekh}

\vspace{3pt}

\small ECE 501 -- Digital Image Processing Course Project

\vspace{2pt}

\small\textbf{Faculty:} Prof. Mehul Raval
\quad $\cdot$ \quad
\textbf{TAs:} Apoorva Mehta \quad $\cdot$ \quad Jyoti Triklani \quad $\cdot$ \quad Sagar Chavda

\vspace{3pt}

\small Date: September 19, 2026

\end{center}

\vspace{1pt}
\hrule height 0.8pt
\vspace{3pt}

% ============================================================
% 1. PROBLEM STATEMENT
% ============================================================

\section{Problem Statement}

This project focuses on the improvement, segmentation, and morphological classification of galaxies captured by images from Hubble Space Telescope (HST) and James Webb Space Telescope (JWST) using classic digital image processing methods. The main difficulty with astronomical images is the presence of noise, changes in illumination or background level, instrument effects, and proximity of astronomical objects to each other.

The purpose of the project is to design a deterministic pipeline that will be able to analyze raw or preliminary processed images from the telescope, detect galaxies in the images, and classify them based on their morphology, such as elliptical, spiral, and irregular galaxies. This project deliberately limits itself to using only classic digital image processing methods instead of machine learning.

% ============================================================
% 2. WEEK 2 FOCUS
% ============================================================

\section{Week 2 Focus}

Week 2 concentrated on progressing from the problem definition set forth in Week 1 to an understanding of the existing
astronomical data and the classical techniques of image processing that may be employed in galaxy morphology studies. The key issues to be researched included:

\begin{itemize}
    \item The classical quantitative metrics of galaxy morphology characterization.
    \item An exploratory study of the image data of HST and JWST.
    \item Selection of appropriate publicly accessible datasets and target galaxies.
    \item Research into segmentation and restoration techniques applicable to astronomical images.
    \item Matching of the selected astronomical targets with publicly available morphological labels from Galaxy Zoo.
\end{itemize}

% ============================================================
% 3. LITERATURE SEARCH AND MORPHOLOGY
% ============================================================

\section{Literature Search on Galaxy Morphology}

The review of literature was focused on classical quantitative characteristics that could be used for discrimination between different types of galaxy morphologies without applying machine learning techniques.
The following set of measurements turned out to be suitable for the project:

\begin{itemize}
    \item \textbf{Concentration:} Concentration – characteristic of the ratio between the brightness of the inner and outer regions of the galaxy.
    \item \textbf{Asymmetry:} Asymmetry – characteristic of the similarity of the galaxy after its rotation or reflection.
    \item \textbf{Smoothness:} – characteristic of small-scale fluctuations of the galaxy structure.
    \item \textbf {Shape characteristics:} – geometrical characteristics such as area, perimeter, eccentricity, and moments could provide additional information about the galaxy morphology. Geometric measurements such as area, perimeter, eccentricity, and moments can provide additional information about morphology.
    \item \textbf {Texture characteristics:} – spatial intensity patterns could be used to distinguish structural features of different galaxies.
\end{itemize}

These descriptors provide a potential basis for constructing interpretable rule-based decision boundaries for the final morphological classification stage.

% ============================================================
% 4. DATASET INVESTIGATION
% ============================================================

\section{Dataset Investigation and Exploratory Data Analysis}

Both HST and JWST imagery are considered during this stage so that the two data sources can eventually be compared using the same classical processing pipeline.

MAST archive (Mikulski Archive for Space Telescopes) is identified as the main source of publicly available
FITS data from HST and JWST telescopes.
The Week 2 study deals with analyzing the characteristics that would influence further stages of image processing:

\begin{itemize}
    \item \textbf{Dynamic Range:} Identification of the intensity range of the astronomical images.
    \item \textbf{Background Levels:} Analysis of background illumination and variation between different images.
    \item \textbf{Noise Characteristics:} Detection of noise in the collected telescope images.
    \item \textbf{Image Resolution:} Comparison of the spatial resolution of the images of HST and JWST.
    \item \textbf{Instrumental Artifacts:} Detection of any artifacts that could hamper segmentation and morphological measurements.
    \item \textbf{Galaxy Separability:} Checking how the galaxies can be distinguished from each other to allow feature extraction.
\end{itemize}

Images of both HST and JWST are analyzed at this stage in order to enable further comparison between two types of data using one classical processing pipeline.

% ============================================================
% 5. DATA AND VALIDATION SOURCES
% ============================================================

\section{Data and Validation Sources}

There are two key sources that have been considered for use in this assignment:
\begin{itemize}
    \item \textbf{NASA MAST}: This is the primary source of astronomical images, especially in FITS format, taken through the HST and JWST telescopes.
    \item \textbf{Galaxy Zoo}: This is the source for publicly available classifications done by citizen scientists and can be used for reference for galaxies

\end{itemize}.
\\
Validation process: Finding common objects between the telescope data and classifications available on Galaxy Zoo.
% ============================================================
% 6. ENVIRONMENT AND TOOLCHAIN
% ============================================================

\section{Environment and Toolchain}

The choice of the computing environment for the project includes Python programming language and traditional image processing libraries.
\begin{itemize}
    \item \textbf{Astropy}: FITS file management and processing of astronomical data.
    \item \textbf{NumPy}: Numerical calculations and manipulation with image arrays.
    \item \textbf{SciPy}: Scientific calculations and signal processing.
    \item \textbf{OpenCV}: Traditional image processing including filtering, thresholding, transformation and morphology.

\end{itemize} 

These libraries contain the computing resources needed for all aspects of classical image processing, from the acquisition of the astronomical images through pre-processing to segmentation, feature extraction, and morphological classification.

% ============================================================
% 7. PRELIMINARY ALGORITHM DIRECTION
% ============================================================

\section{Preliminary Algorithm Direction}

Based on the results of the literature review and data set analysis, the following pipeline will be developed for project implementation:
\begin{enumerate}
    \item \textbf{Image Acquisition} – obtaining sufficient number of HST and JWST FITS images from MAST.
    
    \item \textbf{Image Restoration} – researching classic methods of handling the problem of instrumental blurring and noise, including PSF restoration and deconvolution.
    
    \item \textbf{Background and Contrast} – decreasing the background variability and enhancing galactic objects visibility.
    
    \item \textbf{Galaxy Segmentation} – implementing deterministic filtering, thresholding and morphological processing for separation of individual galaxies.
    
    \item \textbf{Features Extraction} – calculating numerical features for shape, intensity, concentration, asymmetry, smoothness and texture of galaxies.
    
    \item \textbf{Morphological Classification} – developing interpretable classification rules for galaxy morphological classes.

    \item \textbf{Validation} – comparing obtained morphological classifications with Galaxy Zoo morphological classifications.
\end{enumerate}



% ============================================================
% 8. WEEK 2 SUMMARY / STATUS
% ============================================================

\section{Week 2 Summary and Status}

\begin{itemize}
    \item Examined classical quantitative measurements related to the morphology of galaxies, such as concentration, asymmetry, smoothness, shape, and texture.
    \item Chose NASA MAST as the source of the HST and JWST FITS images.
    \item Chose Galaxy Zoo as the external source for verification of the galaxy morphology.
    \item Defined the parameters that are to be investigated within the framework of exploratory data analysis, such as dynamic range, noise, background, resolution, and artifacts of the particular instrument.
    \item Configured the planned Python working environment with Astropy, NumPy, SciPy, and OpenCV libraries.
    \item Developed the initial scheme of the deterministic pipeline of image processing from image acquisition through segmentation, feature extraction, classification, and validation.
    \item \textbf{Status:} The project is currently at the stage when it evolves from problem statement into dataset investigation and literature review.
\end{itemize}


% ============================================================
% 9. OBJECTIVES FOR WEEK 3
% ============================================================

\section{Objectives for Week 3}

\begin{itemize}
    \item \textbf{Dataset Acquisition:} Collect relevant FITS images from the HST and JWST datasets that have adequate galaxies.
    \item \textbf{Initial Image Inspection:} View the obtained FITS images and analyze their intensities and spatial features.
    \item \textbf{Pre-processing:} Find relevant techniques to filter and enhance the images using classical methods.
    \item \textbf{Segmentation Investigation:} Experiment with the thresholding and morphological methods to segment the
galaxies.
    \item \textbf{Target Selection:} ind suitable galaxies that can be classified according to the classifications from Galaxy
Zoo.
    \item \textbf{Pipeline Development:} Start implementing the early phases of the DIP pipeline using classical techniques
with Python.
\end{itemize}



% ============================================================
% END
% ============================================================

\end{document}